% Revised conclusion for ELAPSE (Physica A)

\section{Conclusion}
\label{sec:conclusion}

This paper establishes a rigorous mathematical framework for entropy-triggered
data deletion in distributed networks. The central finding is that the
time-integrated entropy exposure (IEE)---a network-level privacy risk
functional---exhibits a percolation-like phase transition at a critical
entropy threshold $H_c^*$ controlled by the algebraic connectivity $\lambda_2$
of the network Laplacian. Below $H_c^*$, deletion mechanisms contain
propagation efficiently; above $H_c^*$, the IEE diverges as a power law in
network size, with exponent $\alpha \approx 1.55$ consistent across synthetic
topologies. This scaling law, predicted analytically by Corollary~\ref{cor:spectral_decay}
and confirmed numerically for $n \in [50, 500]$, gives network operators a
principled design rule: target $H_c < H_c^*$ to guarantee sub-linear IEE growth.

Among the five mechanisms studied, M5 (social cascade pressure) achieves the
lowest IEE in topology-matched experiments, owing to its sensitivity to local
degree heterogeneity. However, the M6 ensemble---whose weights are learned via
entropy-regularised Nelder--Mead optimisation---provides a robust alternative
under topology uncertainty, network evolution, and threshold misspecification.
The key insight is that ensemble weights aggregated across mechanism types
act as a form of meta-regularisation: no single mechanism dominates unless
it is optimally matched to the unknown true topology, making M6 the safer
default in heterogeneous or unknown-topology deployments.

The theoretical contributions of this work connect ELAPSE to three areas of
statistical physics. First, the phase transition in IEE resembles bond
percolation on random graphs~\cite{Erdos1960}: the critical threshold
$H_c^*(\lambda_2)$ plays the role of the bond occupation probability $p_c$.
Second, the maximum-entropy principle underlying the IEE objective reflects
Jaynes's entropy maximisation framework, casting data deletion as the
\emph{inverse} of entropy maximisation: we seek the control $\Delta(t;\mathbf{w})$
that \emph{minimises} entropy exposure. Third, the spectral trigger time
$\tau(\lambda_2)$---the time at which $\|\mathbf{p}(t) - \mathbf{1}/n\|_2$ crosses a
threshold---connects ELAPSE to the spectral theory of Markov chains and
diffusion on graphs, providing a mechanistic explanation for the $\lambda_2$
dependence observed empirically.

Several directions remain open. The first-passage distribution of entropy
$H(t)$ for the full SDE system is analytically
intractable; a Laplace-method approximation or large-deviation bound
may replace the Gronwall moment bound of Theorem~\ref{thm:iee_bound} in
future work.
Second, the ensemble weight-learning problem is currently solved offline
via Nelder--Mead; online adaptive learning---updating $\mathbf{w}$ as
the network evolves---is an open stochastic control problem.
Third, the dependence of the IEE scaling exponent $\alpha \approx 1.55$ on
the degree exponent $\gamma$ of scale-free networks remains to be characterised
analytically.
Finally, extending ELAPSE to directed or weighted graphs requires replacing
the symmetric Laplacian with a non-self-adjoint operator, for which the
spectral gap theory requires significant modification.

Beyond network privacy, the ELAPSE framework applies to any system where
a diffusing quantity must be contained below a threshold: epidemic control
(data = infection, $H_c$ = herd-immunity threshold), social cascade
suppression (data = misinformation, $H_c$ = virality threshold), and GDPR
compliance (data = personal information copies, $H_c$ = regulatory limit).
In each case, the mathematical structure---SDE with entropy-triggered
mortality, ensemble voting, Nelder--Mead weight learning---transfers
directly, suggesting ELAPSE as a unifying framework for entropy-aware
containment on complex networks.
